\documentclass[11pt]{article}
\usepackage{mathblatt}
% Präambel des Setzers (werkzeuge/setzer-praeambel.tex, wortgleich); dazu \gteil (Muster T2U)
\makeatletter
\setlength{\mbpfbe}{0mm}% keine Punkte (pruefung.md § 5: nur beim Prüfstein)
\newlength{\gEin}\setlength{\gEin}{0pt}
\let\gSchrift\relax
\renewcommand{\pfkreuzzeile}[1]{\par\noindent\hangindent3.2em\hangafter1\hspace*{1.6em}$\square$\ #1\par}
\newcommand{\gkreuz}[1]{$\square$\ #1\hspace{2em}}
\newcommand{\gkreuzzeile}[2]{\par\noindent\hangindent2.2em\hangafter1\mbox{#1\,$\square$}\ #2\par}% Buchstabe, Option: eine je Zeile (neben Skizze)
% links, was man ansieht, rechts, was man tut; Spaltengrenze überall an derselben Stelle (gemeinsam.md „Satz“)
\newsavebox{\gbox}\newlength{\gLinks}
% v0.10 (Lehrer 13.10.): keine starre Halbierung – die linke Spalte ist so breit wie die Figur (+ 6 mm), die
% Antwort steht direkt daneben; ist die Figur breiter als 0,62 Zeile, steht die Antwort darunter
\newcommand{\gzwei}[2]{\par\smallskip\noindent\sbox{\gbox}{#1}%
  \setlength{\gLinks}{\dimexpr\wd\gbox+6mm\relax}%
  \ifdim\gLinks>0.62\linewidth
    \usebox{\gbox}\par\smallskip\noindent\begin{minipage}{\linewidth}#2\end{minipage}%
  \else
    \begin{minipage}[t]{\gLinks}\vspace{0pt}\usebox{\gbox}\end{minipage}%
    \begin{minipage}[t]{\dimexpr\linewidth-\gLinks\relax}\vspace{0pt}\raggedright #2\end{minipage}%
  \fi\par}
% v0.9 (12.10., pruefung.md § 5 „Marke mit B, Bedingungen rechts“): drittes Argument von pfaufg = Bedingungen
% der Prüfung (v0.10: \pfohnetr = TR durchgestrichen, „FOR“), klein grau im rechten Rand auf der Grundlinie der Nummer, Zeilen mit \\
\newcommand{\gbedrand}[2]{\ifblank{#2}{}{\rlap{\hspace*{\dimexpr#1+2mm\relax}%
  \parbox[t]{15mm}{\raggedright\scriptsize\leavevmode\color{mbgrau}\strut#2}}}}
% Teil eines Bündels (a, b …): Marke links im Rand, Bedingungen rechts im Rand, auf der Zeile des Teils
\newcommand{\gteilrand}[2]{\llap{\mbpfrandmarke{#1}\hspace*{\mbpfnr}}\gbedrand{\linewidth}{#2}}
% Teil eines Bündels: Buchstabe in eigener Spalte, Inhalt eingerückt (\gEin = Einzug, damit die Spaltengrenze von
% \gzwei an derselben Stelle bleibt)
\newenvironment{gbteil}[3]{\par\smallskip\noindent\gteilrand{#2}{#3}\makebox[\mbpfteil][l]{\textbf{#1}}%
  \begin{minipage}[t]{\dimexpr\linewidth-\mbpfteil\relax}\setlength{\gEin}{\mbpfteil}\raggedright\strut\ignorespaces}%
 {\end{minipage}\par}
% Grafik höchstens so breit wie die Spalte (Bündel nebeneinander, Koordinatensystem ohne Karo)
\newsavebox{\gpbox}
\newcommand{\gpassend}[1]{\sbox{\gpbox}{#1}\ifdim\wd\gpbox>\linewidth\typeout{GPASSEND: verkleinert \the\wd\gpbox\space auf \the\linewidth}\resizebox{\linewidth}{!}{\usebox{\gpbox}}\else\usebox{\gpbox}\fi}
% Skizzenfeld (v0.10, Lehrer 13.10.): Karo ab dem linken Rand über die Zeile, „Skizze“ klein grau oben links im Feld
% (wie \pfskizzenfeld der Vorlage); der rechte Rand bleibt den Bedingungen
\newcommand{\gskizze}[1]{\par\smallskip\noindent\begin{tikzpicture}
  \clip (0,0) rectangle ({\linewidth-1mm},{#1*10mm});
  \draw[black!16,line width=0.3pt,step=5mm] (0,0) grid ({\linewidth},{#1*10mm});
  \node[anchor=north west,inner sep=2pt,font=\scriptsize,text=mbgrau,fill=white] at (0,{#1*10mm}) {Skizze};
  \end{tikzpicture}\par}
\RenewDocumentEnvironment{pfaufg}{m m m}{%
  \begin{lrbox}{\mbaufgabebox}\begin{minipage}[t]{\linewidth}%
  \gSchrift\noindent\hspace*{\gEin}\mbpfrandmarke{#1}\gbedrand{\linewidth-\gEin}{#3}%
  \begin{minipage}[t]{\mbpfnr}\strut\textbf{#2}\end{minipage}%
  \begin{minipage}[t]{\dimexpr\linewidth-\gEin-\mbpfnr\relax}\raggedright\strut\ignorespaces}%
 {\end{minipage}%
  \end{minipage}\end{lrbox}%
  \setlength{\mbaufgabehoehe}{\dimexpr\ht\mbaufgabebox+\dp\mbaufgabebox\relax}%
  \ifdefempty{\gkopfnach}{\par\addvspace{5pt}\Needspace*{\mbaufgabehoehe}}{%
    % Stufenkopf, der mit dieser (ersten) Aufgabe zusammenbleibt: Platz für beide verlangen, dann den Kopf setzen
    % (12.10., Quadratische: ein Kopf stand allein unten, die Aufgabe mit Koordinatensystem auf der nächsten Seite)
    \par\Needspace*{\dimexpr\mbaufgabehoehe+40pt\relax}\gkopfnach\gdef\gkopfnach{}\par\addvspace{5pt}}%
  \noindent\usebox{\mbaufgabebox}\par\addvspace{5pt}}
% Stufenkopf vormerken: \gkopfvor{<Kopf>} – die nächste pfaufg setzt ihn, sobald beide zusammen Platz haben
\def\gkopfnach{}
\newcommand{\gkopfvor}[1]{\gdef\gkopfnach{#1}}
% Teilaufgabe mit Skizze links in gzwei: Buchstabe in eigener Spalte, darunter Text und Skizze
\newcommand{\gteil}[2]{\makebox[\mbpfteil][l]{\textbf{#1}}\parbox[t]{\dimexpr0.5\textwidth-\gEin-\mbpfnr-5mm-\mbpfteil\relax}{\raggedright #2}}
% Stufenkopf der Vorlage (mathblatt.sty 2026-10-12: Linie, Abstand, 13 pt, \pfjahre) ohne ihren Vorhalt von
% 14 Zeilen: der Setzer fordert vor dem Kopf 0,16 Seitenhöhe (Kopf + erste Aufgabe), mehr nicht
\setlength{\mbpfstufevorhalt}{0pt}
\makeatother
\newcommand{\karo}[1]{\par\smallskip\noindent\begin{tikzpicture}
  \clip (0,0) rectangle ({\linewidth-1mm},{#1*10mm});
  \draw[black!16,line width=0.3pt,step=5mm] (0,0) grid ({\linewidth},{#1*10mm});
  \end{tikzpicture}\par}
\newcommand{\kk}[1]{$\square$\,#1\hspace{0.9em}}
\newcommand{\linie}[1]{\,{\color{black!45}\rule[-1pt]{#1}{0.4pt}}\,}
\newcommand{\pkt}{\textperiodcentered{}}
% v0.7 (12.10.): Schülerfrage unter dem Stufenkopf; Ankreuzoption als Bild (Buchstabe und Kästchen darunter)
\newcommand{\gfrage}[1]{\par\nobreak\noindent\hspace*{\mbpfnr}{\small\itshape #1}\par\nobreak\addvspace{2pt}}
\newlength{\gbildbreite}\setlength{\gbildbreite}{0.32\linewidth}% Setzer setzt sie je Reihe (vier Bilder: 0,24)
\newcommand{\gbildoption}[3]{\begin{minipage}[t][#3][b]{\gbildbreite}\centering #2\par\smallskip\mbox{#1\,$\square$}\end{minipage}}
% Lineare (12.10.): Bild mit Buchstabe ohne Kästchen (Diagramm zuordnen); Aussage mit „wahr / falsch“ zum Ankreuzen
\newcommand{\gbildname}[3]{\begin{minipage}[t][#3][b]{\gbildbreite}\centering #2\par\smallskip\mbox{\textbf{#1}}\end{minipage}}
\newcommand{\gwahrfalsch}[2]{\par\noindent\hangindent2.2em\hangafter1\makebox[2.2em][l]{(#1)}#2\hfill\mbox{wahr $\square$\quad falsch $\square$}\par}
\begin{document}

\pfheftstil
\pflaufkopf[P10 \pkt{} Prüfung]{Wahrscheinlichkeit \pkt{} Lösungen}{VRK2}\pfseitevon
\pfheftkopf{Wahrscheinlichkeit}{Lösungen}
\pfstufenzeile{Ergebnisse aufzählen \pkt{} Wahrscheinlichkeit gesucht \pkt{} Baum ergänzen \pkt{} Pfadregeln anwenden \pkt{} ohne Zurücklegen rechnen \pkt{} Zufallsgerät entwerfen}

\Needspace*{0.16\textheight}\pfstufekopf{Ergebnisse aufzählen}{\pfjahre{25}{20 \pkt{} 17}}
\begin{pfloesungb}{}
\lzb{1.}{}{22, 23, 24, 32, 33, 34, 42, 43, 44}{22, 23, 24, 32, 33, 34, 42, 43, 44 (jede linke Ziffer mit jeder rechten)}
\end{pfloesungb}

\Needspace*{0.16\textheight}\pfstufekopf{Wahrscheinlichkeit gesucht}{\pfjahre{26 \pkt{} 24}{19 \pkt{} 18 \pkt{} 17 \pkt{} 16 \pkt{} 15 \pkt{} 14}}
\begin{pfloesungb}{}
\lzb{2.}{}{P(2) = 50 \%}{P(2) = $\tfrac{3}{6}$ = $\tfrac{1}{2}$ = 50 \%}
\end{pfloesungb}
\begin{pfloesungb}{}
\lzb{3.}{}{P(Staubsaugen) $\approx$ 16,7 \%}{P(Staubsaugen) = $\tfrac{1}{6}$ $\approx$ 16,7 \%}
\end{pfloesungb}

\Needspace*{0.16\textheight}\pfstufekopf{Baum ergänzen}{\pfjahre{26 \pkt{} 24}{20 \pkt{} 19 \pkt{} 15 \pkt{} 14}}
\begin{pfloesungb}{}
\lzb{4.}{}{$\tfrac{5}{8}$}{Gewinn und Niete ergeben zusammen 1: 1 $-$ $\tfrac{3}{8}$ = $\tfrac{5}{8}$; so ergänzt du an jeder Verzweigung den fehlenden Ast}
\end{pfloesungb}
\begin{pfloesungb}{}
\lzb{5.}{}{gerade 2/6, ungerade $\tfrac{4}{6}$; P(beide ungerade) $\approx$ 0,222}{A ungerade $\tfrac{2}{6}$; B ungerade 4/6, nach ungerade: gerade 2/6, ungerade $\tfrac{4}{6}$; P(beide ungerade) = $\tfrac{2}{6}$ · $\tfrac{4}{6}$ = $\tfrac{8}{36}$ = $\tfrac{2}{9}$ $\approx$ 0,222}
\end{pfloesungb}
\begin{pfloesungb}{}
\lzb{6.}{}{kein Joker 5/6, Joker $\tfrac{1}{6}$; P(Joker, Joker) $\approx$ 0,028}{1. Woche: kein Joker 5/6, Joker $\tfrac{1}{6}$; 2. Woche an beiden Ästen: kein Joker 5/6, Joker $\tfrac{1}{6}$; P(Joker, Joker) = $\tfrac{1}{6}$ · $\tfrac{1}{6}$ = $\tfrac{1}{36}$ $\approx$ 0,028}
\end{pfloesungb}

\Needspace*{0.16\textheight}\pfstufekopf{Pfadregeln anwenden}{\pfjahre{26 \pkt{} 25}{20 \pkt{} 14}}
\begin{pfloesungb}{}
\lzb{7.}{}{a) mal; b) je Pfad mal, dann plus}{Pfadregel: die Äste eines Pfades multiplizieren; Summenregel: die Pfade eines Ereignisses addieren}
\end{pfloesungb}
\begin{pfloesungb}{}
\lzb{8.}{a)}{P(34) = 0,125; falsch}{P(34) = $\tfrac{2}{4}$ · $\tfrac{1}{4}$ = $\tfrac{1}{8}$ = 0,125; falsch, P(24) = $\tfrac{1}{4}$ · $\tfrac{1}{4}$ = $\tfrac{1}{16}$ = 0,0625}
\end{pfloesungb}
\begin{pfloesungb}{}
\lzb{}{b)}{P(24 oder 33) = 18,75 \%}{P(24 oder 33) = $\tfrac{1}{4}$ · $\tfrac{1}{4}$ + $\tfrac{2}{4}$ · $\tfrac{1}{4}$ = $\tfrac{3}{16}$ = 18,75 \%}
\end{pfloesungb}
\begin{pfloesungb}{}
\lzb{9.}{}{Gegenereignis „beide gerade“; P(E) $\approx$ 0,778}{Gegenereignis „beide gerade“; P(E) = 1 $-$ $\tfrac{4}{6}$ · $\tfrac{2}{6}$ = 1 $-$ $\tfrac{8}{36}$ = $\tfrac{7}{9}$ $\approx$ 0,778}
\end{pfloesungb}

\Needspace*{0.16\textheight}\pfstufekopf{ohne Zurücklegen rechnen}{\pfjahre{24}{19 \pkt{} 18}}
\begin{pfloesungb}{}
\lzb{10.}{}{4 Kugeln, davon 2 rote}{es liegt eine Kugel weniger da, die roten sind alle noch drin; der nächste Ast „rot“ ist 2/4, nicht $\tfrac{2}{5}$}
\end{pfloesungb}
\begin{pfloesungb}{}
\lzb{11.}{}{Benjamin zieht den Joker nicht ($\tfrac{5}{6}$), Klara aus 5 Zetteln ($\tfrac{1}{5}$); P(Klara Joker) $\approx$ 0,167}{Benjamin zieht den Joker nicht ($\tfrac{5}{6}$), Klara aus 5 Zetteln ($\tfrac{1}{5}$); P(Klara Joker) = $\tfrac{5}{6}$ · $\tfrac{1}{5}$ = $\tfrac{1}{6}$ $\approx$ 0,167}
\end{pfloesungb}

\Needspace*{0.16\textheight}\pfstufekopf{Zufallsgerät entwerfen}{\pfjahre{26 \pkt{} 25}{19 \pkt{} 18}}
\begin{pfloesungb}{}
\lzb{12.}{}{3 und 1}{16 = 4 · 4, also $\tfrac{3}{16}$ = $\tfrac{3}{4}$ · $\tfrac{1}{4}$: ein Rad drei rote Felder, das andere eines}
\end{pfloesungb}
\begin{pfloesungb}{}
\lzb{13.}{}{z. B. links 4, 4, 2, 3 und rechts 2, 3, 4, 3; P(42) = 12,5 \%}{z. B. links 4, 4, 2, 3 und rechts 2, 3, 4, 3; P(42) = $\tfrac{2}{4}$ · $\tfrac{1}{4}$ = $\tfrac{1}{8}$ = 12,5 \%}
\end{pfloesungb}
\begin{pfloesungb}{}
\lzb{14.}{}{z. B. B neu 1, 1, 3, 4, 5, 6 (eine 3 durch 1 ersetzt); P(Summe 3) $\approx$ 0,167}{z. B. B neu 1, 1, 3, 4, 5, 6 (eine 3 durch 1 ersetzt); P(Summe 3) = $\tfrac{3}{6}$ · $\tfrac{2}{6}$ = $\tfrac{6}{36}$ = $\tfrac{1}{6}$ $\approx$ 0,167}
\end{pfloesungb}
\end{document}